\section{Shifted lists and the fixed contour displacement}\label{c7:shifted}

The marked-bond representation uses post-collision displacement \(d=1/8\). Its translated site lists also approach two imaginary heights whose separation is one. We first prove the stability estimate at that displacement, then prove uniform suppression for lists near the two physical heights with the required proportions.

\begin{proposition}[Stability at displacement \(1/8\)]\label{j:stability-eighth}
For the post-collision gas at \(d=1/8\), the conclusion of Theorem~\ref{j:absolute-stability} holds, with constants uniform in \(T\) and in \(\theta\) in a fixed compact subset of \(\mathbb C\).
\end{proposition}

\begin{proof}
We establish the same matrix conclusions as Lemma~\ref{j:stability-coercivity} for this particular displacement. Keep its normalized coordinate \(F=F_{\rm pair}/2\), even direction \(R_d=(U+L)/2\), and odd directions \(O_{\rm odd}=(O_+-O_-)/2\), \(W_d=(U-L)/2\). In this proof put
\[
\Delta=k/8,\qquad c_\Delta=\cosh\Delta,\qquad s_\Delta=\sinh\Delta,
\qquad u=e^{-k/2},\quad z=e^{-k/4},\quad h=\mathsf h_\sigma(k).
\]
The core \(A_\sigma\), and \(M_0,f,g,l,j_0,X,V,\lambda\), are those of \eqref{j:stability-core}.

\paragraph{Schur complements from the height lists.}
Define \(j_c=\cosh(k/2)/(2\cosh k)\), and let \(M^{so}\) denote the short-plus-outer part of the normalized matrix. Summing those kernels on the signed measures gives
\[
\begin{aligned}
M^{so}_{R_d,(P,M,F)}
&=c_\Delta M_{0,F,(P,M,F)}
+s_\Delta\left(j_c,j_c,-\tfrac34+\tfrac1{4\cosh k}\right),\\
M^{so}_{R_dR_d}
&=f c_\Delta^2+\tfrac g2
+\left(-\tfrac12+\tfrac1{2\cosh k}\right)s_\Delta c_\Delta .
\end{aligned}
\]
For example \(p(d)=f c_\Delta-s_\Delta/2\) and
\(o(1/2+d)=j_0c_\Delta-j_cs_\Delta\); applying these identities to the two half-weight heights of \(R_d\) gives the displayed row and square. For
\[
N_\Delta=\frac{R_d-c_\Delta F}{s_\Delta},
\]
the field cross row and square are therefore
\[
M^{so}_{N_\Delta,(P,M,F)}
=\left(j_c,j_c,-\tfrac34+\tfrac1{4\cosh k}\right),
\qquad
M^{so}_{N_\Delta N_\Delta}=\coth\Delta-g/2.
\]
The even \(H\)-moment of \(N_\Delta\) is
\[
Y=-\sinh(3k/2),
\]
and \(V_{N_\Delta}=-1\), since \(V_{R_d}=e^{-\Delta}=c_\Delta-s_\Delta\). The continued entries \eqref{j:height-continuations} give a contribution \(2\) to the \(N_\Delta,F\) cross in the minus channel and zero to its square:
\[
\frac{2\sinh(2\Delta)-4c_\Delta\sinh\Delta}{s_\Delta^2}=0.
\]
Thus the cross column of \(N_\Delta\) with the positive core is exactly
\[
d_1=
\left(j_c,j_c,-\tfrac34+\tfrac1{4\cosh k}
+2\mathbf1_{\{\sigma=-\}}\right)^t
-hXY+\sigma\lambda V.
\]
Its diagonal before the Schur subtraction is
\(\coth\Delta-g/2-hY^2-\sigma\lambda\). The even second Schur complement is consequently
\begin{equation}\label{j:eighth-even-schur}
\coth\Delta-a,\qquad
a=g/2+hY^2+\sigma\lambda+d_1^tA_\sigma^{-1}d_1.
\end{equation}

The odd list gives the first diagonal
\[
G=g+h\sinh^2 k>0
\]
and cross \(c_\Delta r_0+s_\Delta r_1\), where
\[
r_0=j_0+h\sinh k\sinh(3k/2),\qquad
r_1=-j_c+h\sinh k\cosh(3k/2).
\]
Its field square is
\[
g/2-fs_\Delta^2+
\left(\tfrac12-\tfrac1{2\cosh k}\right)s_\Delta c_\Delta.
\]
The uncontinued height square uses
\(h\sinh^2(3k/2+\Delta)\), and its continued correction in the minus channel is
\(-4s_\Delta c_\Delta\). Subtracting the cross square divided by \(G\), and then dividing by \(c_\Delta^2\), gives
\begin{equation}\label{j:eighth-odd-schur}
o_0+b_1\tanh\Delta+b_2\tanh^2\Delta,
\end{equation}
where
\[
\begin{aligned}
o_0&=g/2+h\sinh^2(3k/2)-r_0^2/G,\\
b_2&=-f-g/2+h\cosh^2(3k/2)-r_1^2/G,\\
b_1&=\tfrac12-\tfrac1{2\cosh k}-4\mathbf1_{\{\sigma=-\}}
+2h\sinh(3k/2)\cosh(3k/2)-2r_0r_1/G.
\end{aligned}
\]
Here \(o_0>0\) is the odd Schur complement at zero displacement, whose \(\min(k,1)\) gap follows from the odd Green form in Section~\ref{c7:stability}. These derivations attach both complements to the physical signed lists.

\paragraph{Exact signs and endpoint gaps.}
Substituting the rational core formulas into \eqref{j:eighth-even-schur} and clearing denominators gives
\[
a-1=\frac{u^2\mathsf P_\sigma(u)}
{(1-u^2)\mathsf Q_\sigma(u)}.
\]
The ascending coefficients are
\[
\begin{array}{l|l}
\mathsf P_+&45,12,-13,0,5,-12,29,0,105,-12,-19,0,-37,12,35,0,42\\
\mathsf Q_+&3,36,-44,60,-68,84,-92,108,-78,108,-92,84,-68,60,-44,36,3\\
\mathsf P_-&19,-12,-63,96,31,12,3,48,-177,108,15,0,49,84,-147,48,14\\
\mathsf Q_-&1,12,-40,36,4,12,-16,36,-58,36,-16,12,4,36,-40,12,1 .
\end{array}
\]
Since \(u=z^2\) and \(\tanh\Delta=(1-z)/(1+z)\), define
\[
\mathsf E_\sigma(z)
=2(1+z+z^2+z^3)\mathsf Q_\sigma(z^2)
-z^3\mathsf P_\sigma(z^2).
\]
The sign needed in \eqref{j:eighth-even-schur} is exactly
\[
\coth\Delta-a
=\frac{z\mathsf E_\sigma(z)}
{(1-z^4)\mathsf Q_\sigma(z^2)}.
\]
For the odd complement, substitution into \eqref{j:eighth-odd-schur} gives
\[
\frac{o_0+b_1\tanh\Delta+b_2\tanh^2\Delta}{o_0}
=\frac{4z\,\mathsf T_\sigma(z)}
{(1+z)^4(z^2+1)^2\mathsf D_\sigma(z^4)}.
\]
The ascending first twelve coefficients of the reciprocal degree-\(22\) polynomials \(\mathsf T_+,\mathsf T_-\) are respectively
\[
(4,6,8,8,16,18,20,20,30,33,36,33),\qquad
(4,14,24,24,8,-18,-44,-52,-26,13,52,73),
\]
and the coefficients of \(\mathsf D_+,\mathsf D_-\) are
\[
(2,4,7,4,2),\qquad (10,-16,23,-16,10).
\]
The first denominator is positive by its coefficients, and
\[
\mathsf D_-(y)=(1-y)^2(10+4y+10y^2)+11y^2>0
\quad(0\le y\le1).
\]

For each of the polynomials
\(\mathsf Q_\sigma(z^2),\mathsf E_\sigma(z),\mathsf T_\sigma(z)\),
write its ascending coefficients as \(p_i\) and its degree as \(s\). The coefficient of \(r^j\) in
\((1+r)^sp(1/(1+r))\) is
\[
\sum_{i=0}^{s-j}\binom{s-i}{j}p_i.
\]
Direct binomial expansion gives the following minimum coefficient in each case:
\[
\begin{array}{c|c|cc}
 &s&\sigma=+&\sigma=-\\\hline
\mathsf Q_\sigma(z^2)&32&3&1\\
\mathsf E_\sigma(z)&35&6&2\\
\mathsf T_\sigma(z)&22&4&4 .
\end{array}
\]
All coefficients are positive. Taking \(r=(1-z)/z\) proves positivity for \(0<z\le1\), and the constant coefficients give positivity at \(z=0\). This proves both Schur signs for every \(k>0\).

The same identities give the following limits as \(k\downarrow0\):
\[
\begin{array}{c|cc}
 &\sigma=+&\sigma=-\\\hline
k(\coth\Delta-a)&6&4\\
k^{-1}\operatorname{Schur}_{R_d-c_\Delta F}&3/32&1/16\\
G/k&11/18&3/2\\
k^{-1}\operatorname{Schur}_{W_d}&431/1408&71/384 .
\end{array}
\]
Here the second row restores the factor \(s_\Delta^2\), and the last row restores \(c_\Delta^2\) in the odd complement. The cross column of \(R_d-c_\Delta F\) is \(s_\Delta d_1=O(k)\). Its total charge is \(\cosh\Delta-1=O(k^2)\); also \(q^tA_+^{-1}=O(k)\) by the core expansion. Completing the square with the core therefore retains a positive part of the common-charge bound. In the odd block all entries are \(O(k)\), so its two positive pivot limits give a \(k\)-order gap.

For high frequencies return to the fixed basis \(P,M,F,R_d\) and \(O_{\rm odd},W_d\). The \(R_d\) crosses decay to zero and its extra diagonal tends to \(1/2\), as does the extra odd diagonal. In the minus channel the growing extreme terms cancel by \eqref{j:height-continuations}; their remainders decay because the next height rate is \(4\) and \(3+2d=13/4<4\). Thus the fixed-basis matrices converge exponentially to positive quotient limits. The count null relation and the entrywise nonnegative physical limit are unchanged. The construction of \(J\) and the Fourier identity \eqref{j:stability-energy} now apply at \(d=1/8\), proving the proposition.
\end{proof}

\subsection{Suppression for translated site lists}

For one species, a site at imaginary height \(h\) contributes to a suppressed particle \(\tau\) the factor
\[
\omega_\tau(x;h)=B(x+i(a_\tau-h))^{\epsilon_\tau},
\qquad
\begin{array}{c|ccccc}
\tau&P&O_+&O_-&U&L\\\hline
a_\tau&0&-1&1&13/8&-13/8\\
\epsilon_\tau&1&1&1&-1&-1 .
\end{array}
\]
The \(M\) and \(F_{\rm pair}\) site factors are exactly one after the per-site meromorphic cancellations supplied by the identities for \(B\). These are statements about the \(B\)-products; the additional marked-source dresses are separate.

\begin{proposition}[Uniform suppression near the two physical heights]\label{j:shifted-suppression}
There are constants \(0<\delta_0<1/16\) and \(c_1>0\) with the following property. For every fixed discrepancy bound \(0\le C_0<\infty\) there is \(n_0(C_0)\) such that, for integers \(n\ge n_0(C_0)\), the conclusion below holds uniformly in \(f\in[0,3/5]\).

Partition a list of \(n\) real heights into groups \(I_0,I_1\) such that
\[
\bigl||I_0|-n(1-f)\bigr|\le C_0,\qquad
\bigl||I_1|-nf\bigr|\le C_0.
\]
Suppose
\[
|h_a+f|<\delta_0\quad(a\in I_0),\qquad
|h_a-(1-f)|<\delta_0\quad(a\in I_1).
\]
For every real \(x\) and every \(\tau\in\{P,O_+,O_-,U,L\}\),
\begin{equation}\label{j:shifted-suppression-bound}
\left|\prod_{a=1}^n\omega_\tau(x;h_a)\right|
\le \exp\!\left(-c_1 n e^{-\kappa_1|x|}\right),
\qquad \kappa_1=\pi/3.
\end{equation}
The sums of the one-site logarithms chosen to vanish at either infinity are \(O(n e^{-\kappa_1|x|})\) there, uniformly in these lists. The proposition applies separately to each species.
\end{proposition}

\begin{proof}
We first prove strict suppression at the exact heights, with proportions \(1-f,f\). The \(U,L\) factors are the delicate ones. For real \(x\), put
\[
C_1=\cosh(\kappa_1x),\qquad
s_0=\frac{(\sqrt3/2)C_1}{\sqrt{C_1^2-1/4}},\qquad
R=\frac{\sin(\kappa_1t)}{\sqrt{C_1^2-1/4}}.
\]
Then
\[
b(t):=\log|B(x+i(3/2+t))|
=\operatorname{atanh}\frac{2s_0R}{1+R^2}.
\]
Unless both terms already have the favorable sign, the required comparison for \(U\) is
\[
(1-f)b(f+1/8)>f b(7/8-f),\qquad 0\le f\le7/8.
\]
The \(L\) comparison is the same after using \(1-f\) in place of \(f\); outside the displayed mixed-sign range both terms are favorable.

For \(0\le t\le1\), with limiting values at the endpoints,
\[
\frac9{10}\le\frac{b(t)}{2s_0R}
\le1+\frac{\sin^2(\kappa_1t)}
{3(3/4-\sin^2(\kappa_1t))}.
\]
For the upper bound, \(s_0\le1\) reduces the ratio to at most
\(\operatorname{atanh}(R)/R\). Its positive series is at most
\(1+R^2/(3(1-R^2))\), and \(R^2\le\sin^2(\kappa_1t)/(3/4)\) gives the displayed bound. For the lower bound, use
\(s_0^2\ge(3+R^2)/4\) and
\[
v\coth v\ge1+v^2/3-v^4/45.
\]
The latter follows by multiplying by \(\sinh v\) and comparing its power series. The desired lower bound is equivalent to \(v\coth v\ge(9/10)(1+R^2)\) at \(v=(9/5)s_0R\). The polynomial on the right increases with \(v^2\) for \(0\le v^2\le81/25\), so we may substitute the lower bound for \(s_0^2\). The difference then has numerator
\[
5000-4500y+6939y^2-4374y^3-729y^4,\qquad y=R^2\in[0,1],
\]
over \(50000\). Under \(y=r/(1+r)\), multiplication by \((1+r)^4\) gives the positive coefficient list
\[
(5000,15500,23439,16004,2336).
\]
This proves the lower ratio bound.

Put \(t=7/8-f\). Concavity of sine gives
\[
\sin(\kappa_1(f+1/8))\ge(\sqrt3/2)(1-t),
\]
and its Taylor estimate, using \(\kappa_1<21/20\), gives
\[
\sin(\kappa_1t)\le p(t):=(21/20)t-(9/50)t^3<\sqrt3/2.
\]
The preceding ratio bounds reduce the strict comparison to
\[
\frac{77}{100}(1/8+t)(1-t)3(3/4-p(t)^2)
-(7/8-t)p(t)(9/4-2p(t)^2)>0,\qquad 0\le t\le7/8.
\]
After \(t=(7/8)r/(1+r)\) and multiplication by \((1+r)^{10}\), direct binomial expansion gives every coefficient greater than \(1/50\). This proves the \(U,L\) comparison. For \(O_\pm\), the mixed-sign comparison is
\[
(1-f)b(f+1/2)>f b(1/2-f),\qquad 0\le f\le1/2,
\]
or its exchanged version, and follows from monotonicity of \(b\) on \([0,1]\). The \(P\) factor is strictly suppressed because both exact heights lie in \(|\Im v|<3/2\).

We make the uniform neighborhood in the proposition explicit. By \(B(v+3i)=1/B(v)\), all five one-site factors can also be written
\[
\omega_\tau(x;h)=B(x+i(\widetilde a_\tau-h)),
\qquad
(\widetilde a_P,\widetilde a_{O_+},\widetilde a_{O_-},
\widetilde a_U,\widetilde a_L)=(0,-1,1,-11/8,11/8).
\]
If a height is within \(1/16\) of either \(-f\) or \(1-f\), then every imaginary argument here lies in
\[
[-39/16,\,163/80]\subset(-5/2,5/2).
\]
On this tube the factors have no poles; their only possible singular behavior in the negative logarithm is \(+\infty\) at a zero.

Define the rescaled energy
\[
\Phi_\tau(x,h)
=-e^{\kappa_1|x|}\log|B(x+i(\widetilde a_\tau-h))|,
\]
assigning \(+\infty\) at zeros. The expansion at either infinity gives the continuous limiting value
\[
\Phi_\tau(\pm\infty,h)=2\sqrt3\cos(\kappa_1(\widetilde a_\tau-h)).
\]
Let \(J_h=[-53/80,17/16]\), which contains every height in the preliminary tube. For each finite \(M\), the truncation \(\Phi_{\tau,M}=\min(\Phi_\tau,M)\) is continuous on the compact space \(\overline{\mathbb R}\times J_h\), where \(\overline{\mathbb R}\) adjoins \(x=\pm\infty\). At infinity the rescaled \(b(t)\) is \(2\sqrt3\sin(\kappa_1t)\), and the same strict ideal comparisons persist. Thus the ideal comparisons say that
\[
(1-f)\Phi_\tau(x,-f)+f\Phi_\tau(x,1-f)>0
\]
at every \((x,f)\in\overline{\mathbb R}\times[0,3/5]\), with the value \(0\cdot(+\infty)=0\). The truncated expressions increase to this positive expression as \(M\to\infty\). At each point some finite \(M\) gives a positive margin; continuity and a finite cover therefore provide one \(M\) and one \(a_0>0\), common to all five labels, such that
\[
(1-f)\Phi_{\tau,M}(x,-f)+f\Phi_{\tau,M}(x,1-f)\ge a_0
\]
uniformly in \(x,f,\tau\).

Uniform continuity of these finitely many truncated functions now gives one \(0<\delta_0<1/16\) such that changing a height by less than \(\delta_0\) changes \(\Phi_{\tau,M}\) by at most \(a_0/4\), uniformly in \(x\) and the target height. Let \(B_0\) bound their absolute values. For the perturbed list in the proposition,
\[
\begin{aligned}
\sum_{a=1}^n\Phi_\tau(x,h_a)
&\ge\sum_{a=1}^n\Phi_{\tau,M}(x,h_a)\\
&\ge n a_0-\tfrac14na_0-2C_0B_0.
\end{aligned}
\]
For \(n\ge n_0(C_0)\) the last expression is at least \(na_0/2\). This proves \eqref{j:shifted-suppression-bound} with \(c_1=a_0/2\), independent of \(f\) and \(n\). The uniform expansion of \(B\) on the same tube gives the asserted logarithmic bounds at infinity. Applying \eqref{j:shifted-suppression-bound} to each image gives the corresponding periodized suppression.
\end{proof}
